% GENERATED FILE. Edit canonical lessons, then run npm run generate:books.
% canonical-source-hash: fnv1a64:5804df3f6beabb2b
% canonical-lessons: GU-C449-atur, GU-C449-utsuk, GU-C449-utkantha, GU-C449-rijhvu, GU-C449-prasann

\chapter{Waiting Eagerly}
\label{ch:gu-waiting-eagerly}
\hlchaptermodality{449}

\begin{chapteropening}
\textbf{By the end of this chapter:} \emph{I can say how eager I am for something.}

Complete the last lesson of chapter 449: I can say how eager I am for something.
\end{chapteropening}

\section[\texorpdfstring{\gu{આતુર} (ātur)}{આતુર (ātur)}]{\gu{આતુર} (ātur) — eager; impatient for}

\label{lesson:GU-C449-atur}



\begin{quote}
Before the new one: say the Gujarati for a joke, then the Gujarati for to tease.
\end{quote}

\subsection*{You'll want to know: \gu{આતુર}}
\textbf{\gu{આતુર}} — \emph{ātur} — \textquotedblleft{}eager; impatient for\textquotedblright{}.

\begin{quote}
\gu{બાળકો} \gu{દિવાળીની} \gu{રજા} \gu{માટે} \gu{આતુર} \gu{છે}.
\end{quote}

\begin{quote}
\emph{bāḷko divāḷīnī rajā māṭe ātur che.}
\end{quote}

\begin{quote}
The children are eager for the Diwali holidays.
\end{quote}

\begin{cousinweb}
\textbf{\gu{આતુર}} (\emph{ātur}) is a tatsama: Monier-Williams gives Sanskrit \emph{ātura} as \textquotedblleft{}suffering, sick in body or mind\textquotedblright{} and also as \textquotedblleft{}desirous of\textquotedblright{}; Gujarati keeps that second sense, eager and restless for something.
\end{cousinweb}

\begin{grammarlens}[title={What you've built}]
One more word for describing what you lived and felt.
\end{grammarlens}

\subsection*{Your turn}
\begin{itemize}
\raggedright
  \item \emph{Say it:} \emph{ātur}
  \item \emph{Say it:} \emph{ātur}, once more
  \item \emph{Say it:} \emph{bāḷko divāḷīnī rajā māṭe ātur che.} — The children are eager for the Diwali holidays.
  \item \emph{Recall:} say the Gujarati for a joke, then the Gujarati for to tease, then say \emph{ātur} again
\end{itemize}

\begin{tcolorbox}[breakable,colback=teal!4,colframe=teal!35!black,title={Before you move on}]
What does \gu{આતુર} mean? (\textquotedblleft{}Eager; impatient for\textquotedblright{}.) Say it once more.
\end{tcolorbox}
\section[\texorpdfstring{\gu{ઉત્સુક} (utsuk)}{ઉત્સુક (utsuk)}]{\gu{ઉત્સુક} (utsuk) — keen; curious to know}

\label{lesson:GU-C449-utsuk}



\begin{quote}
Before the new one: say the Gujarati for to tease, then the Gujarati for eager.
\end{quote}

\subsection*{You'll want to know: \gu{ઉત્સુક}}
\textbf{\gu{ઉત્સુક}} — \emph{utsuk} — \textquotedblleft{}keen; curious to know\textquotedblright{}.

\begin{quote}
\gu{મીરા} \gu{પરિણામ} \gu{જાણવા} \gu{માટે} \gu{ઉત્સુક} \gu{હતી}.
\end{quote}

\begin{quote}
\emph{Mīrā pariṇām jāṇvā māṭe utsuk hatī.}
\end{quote}

\begin{quote}
Mira was keen to know the result.
\end{quote}

\begin{cousinweb}
\textbf{\gu{ઉત્સુક}} (\emph{utsuk}) is a tatsama, Sanskrit \emph{utsuka}, \textquotedblleft{}anxiously desirous, striving for\textquotedblright{}, as Monier-Williams records it.
\end{cousinweb}

\begin{grammarlens}[title={What you've built}]
One more word for describing what you lived and felt.
\end{grammarlens}

\subsection*{Your turn}
\begin{itemize}
\raggedright
  \item \emph{Say it:} \emph{utsuk}
  \item \emph{Say it:} \emph{utsuk}, once more
  \item \emph{Say it:} \emph{Mīrā pariṇām jāṇvā māṭe utsuk hatī.} — Mira was keen to know the result.
  \item \emph{Recall:} say the Gujarati for to tease, then the Gujarati for eager, then say \emph{utsuk} again
\end{itemize}

\begin{tcolorbox}[breakable,colback=teal!4,colframe=teal!35!black,title={Before you move on}]
What does \gu{ઉત્સુક} mean? (\textquotedblleft{}Keen; curious to know\textquotedblright{}.) Say it once more.
\end{tcolorbox}
\section[\texorpdfstring{\gu{ઉત્કંઠા} (utkanṭhā)}{ઉત્કંઠા (utkanṭhā)}]{\gu{ઉત્કંઠા} (utkanṭhā) — longing; eager desire}

\label{lesson:GU-C449-utkantha}



\begin{quote}
Before the new one: say the Gujarati for eager, then the Gujarati for keen.
\end{quote}

\subsection*{You'll want to know: \gu{ઉત્કંઠા}}
\textbf{\gu{ઉત્કંઠા}} — \emph{utkanṭhā} — \textquotedblleft{}longing; eager desire\textquotedblright{}.

\begin{quote}
\gu{માતાને} \gu{દીકરીને} \gu{જોવાની} \gu{ઘણી} \gu{ઉત્કંઠા} \gu{હતી}.
\end{quote}

\begin{quote}
\emph{mātāne dīkrīne jovānī ghaṇī utkanṭhā hatī.}
\end{quote}

\begin{quote}
The mother longed to see her daughter.
\end{quote}

A feminine noun.

\begin{cousinweb}
\textbf{\gu{ઉત્કંઠા}} (\emph{utkanṭhā}) is a tatsama, Sanskrit \emph{utkaṇṭhā}, \textquotedblleft{}longing for a beloved person or thing\textquotedblright{}, as Monier-Williams records it.
\end{cousinweb}

\begin{grammarlens}[title={What you've built}]
One more word for describing what you lived and felt.
\end{grammarlens}

\subsection*{Your turn}
\begin{itemize}
\raggedright
  \item \emph{Say it:} \emph{utkanṭhā}
  \item \emph{Say it:} \emph{utkanṭhā}, once more
  \item \emph{Say it:} \emph{mātāne dīkrīne jovānī ghaṇī utkanṭhā hatī.} — The mother longed to see her daughter.
  \item \emph{Recall:} say the Gujarati for eager, then the Gujarati for keen, then say \emph{utkanṭhā} again
\end{itemize}

\begin{tcolorbox}[breakable,colback=teal!4,colframe=teal!35!black,title={Before you move on}]
What does \gu{ઉત્કંઠા} mean? (\textquotedblleft{}Longing; eager desire\textquotedblright{}.) Say it once more.
\end{tcolorbox}
\section[\texorpdfstring{\gu{રીઝવું} (rījhvũ)}{રીઝવું (rījhvũ)}]{\gu{રીઝવું} (rījhvũ) — to be pleased; to be won over}

\label{lesson:GU-C449-rijhvu}



\begin{quote}
Before the new one: say the Gujarati for keen, then the Gujarati for longing.
\end{quote}

\subsection*{You'll want to know: \gu{રીઝવું}}
\textbf{\gu{રીઝવું}} — \emph{rījhvũ} — \textquotedblleft{}to be pleased; to be won over\textquotedblright{}.

\begin{quote}
\gu{દાદા} \gu{નાની} \gu{ભેટથી} \gu{પણ} \gu{રીઝે} \gu{છે}.
\end{quote}

\begin{quote}
\emph{dādā nānī bheṭthī paṇ rījhe che.}
\end{quote}

\begin{quote}
Grandfather is pleased even by a small gift.
\end{quote}

Past (he …): \gu{તે} \gu{રીઝ્યો}, \emph{te rījhyo}.

\begin{cousinweb}
\textbf{\gu{રીઝવું}} (\emph{rījhvũ}) is inherited: CDIAL lists the Gujarati verb, \textquotedblleft{}to be pleased\textquotedblright{}, under Sanskrit \emph{ṛdhati}, \textquotedblleft{}increases, prospers\textquotedblright{} (CDIAL 2457).
\end{cousinweb}

\begin{grammarlens}[title={What you've built}]
One more word for describing what you lived and felt.
\end{grammarlens}

\subsection*{Your turn}
\begin{itemize}
\raggedright
  \item \emph{Say it:} \emph{rījhvũ}
  \item \emph{Say it:} \emph{rījhvũ}, once more
  \item \emph{Say it:} \emph{dādā nānī bheṭthī paṇ rījhe che.} — Grandfather is pleased even by a small gift.
  \item \emph{Recall:} say the Gujarati for keen, then the Gujarati for longing, then say \emph{rījhvũ} again
\end{itemize}

\begin{tcolorbox}[breakable,colback=teal!4,colframe=teal!35!black,title={Before you move on}]
What does \gu{રીઝવું} mean? (\textquotedblleft{}To be pleased; to be won over\textquotedblright{}.) Say it once more.
\end{tcolorbox}
\section[\texorpdfstring{\gu{પ્રસન્ન} (prasann)}{પ્રસન્ન (prasann)}]{\gu{પ્રસન્ન} (prasann) — pleased; cheerful}

\label{lesson:GU-C449-prasann}



\begin{quote}
Before the new one: say the Gujarati for longing, then the Gujarati for to be won over.
\end{quote}

\subsection*{You'll want to know: \gu{પ્રસન્ન}}
\textbf{\gu{પ્રસન્ન}} — \emph{prasann} — \textquotedblleft{}pleased; cheerful\textquotedblright{}.

\begin{quote}
\gu{ગુરુ} \gu{મારા} \gu{જવાબથી} \gu{પ્રસન્ન} \gu{થયા}.
\end{quote}

\begin{quote}
\emph{guru mārā javābthī prasann thayā.}
\end{quote}

\begin{quote}
The teacher was pleased with my answer.
\end{quote}

\begin{cousinweb}
\textbf{\gu{પ્રસન્ન}} (\emph{prasann}) is a tatsama, Sanskrit \emph{prasanna}, \textquotedblleft{}clear, bright, pure\textquotedblright{}, as Monier-Williams records it; Gujarati uses it for a clear, pleased mind.
\end{cousinweb}

\begin{grammarlens}[title={What you've built}]
One more word for describing what you lived and felt.
\end{grammarlens}

\subsection*{Your turn}
\begin{itemize}
\raggedright
  \item \emph{Say it:} \emph{prasann}
  \item \emph{Say it:} \emph{prasann}, once more
  \item \emph{Say it:} \emph{guru mārā javābthī prasann thayā.} — The teacher was pleased with my answer.
  \item \emph{Recall:} say the Gujarati for longing, then the Gujarati for to be won over, then say \emph{prasann} again
\end{itemize}

\begin{tcolorbox}[breakable,colback=teal!4,colframe=teal!35!black,title={Before you move on}]
What does \gu{પ્રસન્ન} mean? (\textquotedblleft{}Pleased; cheerful\textquotedblright{}.) Say it once more.
\end{tcolorbox}
